\section{Related Work}

\subsection{Temperature Scaling and Post-Hoc Calibration}

Temperature scaling, introduced by \cite{guo2017calibration}, optimizes a single parameter $T$ to rescale logits before softmax, improving Expected Calibration Error (ECE) for modern neural networks. This method incurs zero inference overhead since $T$ is applied post-training. However, Guo et al. focused on calibration metrics (ECE), not selective prediction AUROC---our work extends this by evaluating temperature scaling's effectiveness for rejecting incorrect predictions.

More recent work by \cite{nixon2019ensemble} explored ensemble temperature scaling, and \cite{desai2020calibration} studied calibration in NLG tasks. \textbf{None report AUROC vs cost trade-offs}, instead treating calibration and selective prediction as separate problems.

\subsection{Conformal Prediction for Language Models}

Conformal prediction provides distribution-free coverage guarantees for uncertainty sets. \cite{kumar2023conformal} introduced conformal prediction for selective NLG, demonstrating validity under distribution shift. \cite{su2024conformal} extended this to API-only settings (no logit access) using sample frequency and semantic similarity for nonconformity scores, achieving competitive performance on hallucination detection.

While these methods establish conformal prediction's viability for LLMs, \textbf{they do not compare inference cost across multiple UQ approaches}. Our work tests conformal prediction's cross-dataset generalization (HaluEval calibration $\rightarrow$ TruthfulQA test) and positions it within the cost-performance space.

\subsection{Monte Carlo Dropout}

\cite{gal2016dropout} showed that dropout at inference time approximates Bayesian inference, enabling uncertainty estimation via variance across $k$ stochastic forward passes. This approach has been applied to vision tasks \cite{kendall2017uncertainties} and NLP \cite{xiao2019quantifying}, typically using $k=10$-30 samples.

Recent work on LLM uncertainty includes \cite{kuhn2023semantic} on semantic entropy and \cite{lin2023conformal} on conformal language modeling. However, \textbf{$k$-dependency for LLMs remains understudied}---prior work does not systematically vary $k$ to identify efficiency sweet spots. Our finding that $k=3$ suffices for AUROC $\geq$ 0.70 at 8B scale contrasts with vision domain conventions ($k=30$ in retinal-selective-prediction).

\subsection{Selective Prediction Benchmarks}

\cite{yang2023selective} applied selective prediction to TruthfulQA, showing that uncertainty-based rejection improves accuracy. Their work validates TruthfulQA as a selective prediction benchmark but \textbf{does not compare multiple UQ methods or report inference costs}.

\cite{zhang2020calibration} compared ensemble and compositional calibration methods, introducing kernel-density ECE estimators. While comprehensive on calibration metrics, this work predates LLM-scale models and lacks selective prediction evaluation.

\subsection{Our Positioning}

Existing work evaluates UQ methods in isolation or compares methods without cost analysis. \textbf{We are the first to systematically benchmark AUROC vs inference cost} for multiple UQ method variants on the same selective prediction task. Our Pareto frontier framing shifts the question from ``which method wins'' to ``which method for my budget,'' enabling practitioners to make informed trade-offs rather than defaulting to the most expensive option.
